\chapter{Introducción}

\section{Contexto}

BookPoint Chile es una empresa nacional dedicada a la venta de libros, artículos de
papelería y material educativo. Nació como una librería independiente en Concepción y,
debido al crecimiento sostenido de sus ventas, abrió nuevas sucursales en Temuco y La
Serena.

Actualmente, la empresa opera con un sistema monolítico que concentra ventas,
inventario, despacho y atención de clientes en un único aplicativo. El aumento
sostenido de pedidos, tanto presenciales como por canal web, ha comenzado a provocar
lentitud en los procesos, problemas de stock desactualizado entre sucursales y demoras
en la gestión de despachos.

\section{Objetivo del informe}

Este informe documenta la propuesta de solución tecnológica desarrollada para
modernizar el sistema de BookPoint Chile, migrándolo desde su arquitectura monolítica
actual hacia una arquitectura de microservicios independientes. En concreto, cubre los
siete aspectos exigidos por la pauta de la Evaluación Final Transversal para la
Situación Evaluativa 1 (Entrega por Encargo):

\begin{enumerate}
    \item La estrategia de microservicios utilizada para resolver el caso propuesto.
    \item La evaluación de los enfoques éticos considerados en su implementación.
    \item Las buenas prácticas de diseño, arquitectura, y las herramientas (Spring,
    Maven, Git) utilizadas para desarrollar los componentes de backend.
    \item La implementación de operaciones CRUD con JPA/ORM, validadas con Postman.
    \item La comunicación RESTful entre microservicios, validada con Postman.
    \item El desarrollo de pruebas unitarias con JUnit.
    \item La implementación de un framework de documentación HATEOAS.
\end{enumerate}

\section{Alcance del proyecto}

La solución desarrollada consta de \textbf{diez microservicios de negocio}
(catálogo, inventario, sucursales, clientes, carrito, pedidos, pagos, envíos,
notificaciones y reseñas), además de un servidor de descubrimiento (Eureka) y un API
Gateway como punto de entrada único.

Se priorizó cubrir en profundidad el \textbf{flujo comercial completo} descrito en el
caso -- desde que un cliente explora el catálogo hasta que recibe su pedido y puede
dejar una reseña --, con validaciones reales entre servicios, manejo de errores
consistente y pruebas automatizadas. Quedan deliberadamente fuera del alcance actual
funciones administrativas de soporte descritas en el caso (autenticación y roles,
reportes de venta, promociones, facturación electrónica, transferencias entre
sucursales), que se detallan junto con su justificación en la
sección~\ref{ch:arquitectura}.
